\documentclass[12pt]{article}
\usepackage[margin=1in]{geometry}
\usepackage[utf8]{inputenc}
\usepackage{latexsym, amsfonts, enumitem, graphicx, environ,enumitem, fancyhdr, color, amssymb, amsthm, amsmath, mathtools, tikz, nicematrix}
\usepackage{physics}
\usetikzlibrary{patterns}
\usetikzlibrary{decorations.pathreplacing}
\pagestyle{fancy}
\setlength{\headheight}{65pt}
\newenvironment{problem}[2][Problem]{\begin{trivlist}
    \item[\hskip \labelsep {\bfseries #1}\hskip \labelsep {\bfseries #2.}]}{\end{trivlist}}
\DeclareMathOperator{\Ima}{Im}

\usepackage{hyperref}
\usepackage[capitalize]{cleveref}
\usepackage{thmtools}

\newtheorem{thm}{Theorem}[section]
\newtheorem{lem}[thm]{Lemma}
\newtheorem{cl}[thm]{Claim}
\newtheorem{prop}[thm]{Proposition}
\newtheorem{cor}[thm]{Corollary}
\newtheorem{conj}[thm]{Conjecture}
\newtheorem{obstruction}[thm]{Obstruction}

\usepackage{mdframed}

\theoremstyle{definition}
\newtheorem{definition}[thm]{Definition}
\newtheorem{remark}[thm]{Remark}
\newtheorem{question}{Question}
\newtheorem{example}[thm]{Example}
\newtheorem{exercise}[thm]{Exercise}
\newtheorem{properties}[thm]{Properties}
\newtheorem*{answer}{Answer}

\usepackage{tcolorbox}
\tcbuselibrary{theorems,skins,breakable}

\tcolorboxenvironment{example}{
enhanced jigsaw,colframe=cyan,interior hidden,
breakable,%before skip=10pt,after skip=10pt
}

\tcolorboxenvironment{thm}{
enhanced jigsaw,colframe=red,interior hidden,
breakable,%before skip=10pt,after skip=10pt
}

\tcolorboxenvironment{prop}{
enhanced jigsaw,colframe=red,interior hidden,
breakable,%before skip=10pt,after skip=10pt
}

\tcolorboxenvironment{properties}{
enhanced jigsaw,colframe=red,interior hidden,
breakable,%before skip=10pt,after skip=10pt
}

\tcolorboxenvironment{lem}{
enhanced jigsaw,colframe=red,interior hidden,
breakable,%before skip=10pt,after skip=10pt
}
\tcolorboxenvironment{cor}{
enhanced jigsaw,colframe=red,interior hidden,
breakable,%before skip=10pt,after skip=10pt
}
\tcolorboxenvironment{obstruction}{
enhanced jigsaw,colframe=red,interior hidden,
breakable,%before skip=10pt,after skip=10pt
}

\tcolorboxenvironment{proof}{
enhanced jigsaw, frame hidden,%interior hidden,
breakable,%before skip=10pt,after skip=10pt
}

\tcolorboxenvironment{answer}{
enhanced jigsaw, frame hidden,%interior hidden,
breakable,%before skip=10pt,after skip=10pt
}

\renewcommand\qedsymbol{\( \blacksquare \)}

\lhead{Avi Chetlin \\ Assignment 02 \\ 10 September 2026}
\rhead{Dr. Michael Martens \\ PHYS 313 Thermodynamics \\ Fall 2026}

\begin{document}
\begin{problem}{1}
  Multiplicity of some system as a function of energy is given as
  \begin{equation}
  \label{eq:1}
  g(U) = C U^{\frac{3N}{2}},
  \end{equation}
  where \( C \) is a constant and \( N \) is the number of particles.
  The system is in thermal equilibrium with a reservoir at temperature \(\tau\).
  \begin{enumerate}[label=(\alph*)]
    \item Show that the equilibrium energy of the system is \( \hat{U} = \frac{3}{2}N \tau\).
          \begin{proof}
            The equilibrium energy occurs when
            \begin{equation}
            \label{eq:2}
            \frac{\partial \sigma_{S}}{\partial U} = \frac{\partial \sigma_{R}}{\partial U} = \frac{1}{\tau}.
            \end{equation}

            \(\sigma_S = \log g(U) = \log C + \log U^{\frac{3N}{2}} = \log C + \frac{3N}{2}\log U\), so
            \begin{equation}
            \label{eq:3}
            \frac{\partial \sigma_{S}}{\partial U} = \frac{3N}{2U}.
            \end{equation}
            If \( \hat{U} \) is the equilibrium energy, then
            \begin{equation}
            \label{eq:4}
            \frac{3N}{2 \hat{U}} = \frac{1}{\tau},
          \end{equation}
          hence \( \hat{U} = \frac{3N\tau}{2} \), as desired.
          \end{proof}
    \item Show that \( \frac{\partial^2 \sigma}{\partial U^2} \eval_{U = \hat{U}}\) is negative.
          \begin{proof}
            \begin{equation}
            \label{eq:5}
            \frac{\partial^2 \sigma_{S}}{\partial U^2} = \frac{\partial }{\partial U} \frac{3N}{2U} = - \frac{3N}{2 U^2}.
            \end{equation}
            This is always negative, since \( N \) is necessarily positive as is \( U^2 \).
            In particular, evaluated at \( \hat{U} \), \( \frac{\partial^2 \sigma_{S}}{\partial U^2} \eval_{U = \hat{U}} < 0  \).
          \end{proof}
  \end{enumerate}
\end{problem}

\begin{problem}{2}
  Based on the figure found in the problem sheet.
  The two systems initially have the values indicated, and are then brought into thermal contact.
  \begin{enumerate}[label=(\alph*)]
    \item Explain what happens subsequently and why, without using the word temperature.
          What is the result after a long time?
          \begin{answer}
            Initially, as the energy from the system with a shallower slope (higher temperature) flows into the system with a steeper slope (lower temperature), the slope of the graph for system \( A \) will decrease, and the slope for \( B \) will increase.
            This is because systems always evolve with increasing entropy towards equilibrium, which occurs when \(\partial_U\sigma_A = \partial_U\sigma_B\), i.e., when the slopes are equal.

            After a long time, the two slopes will be equal, and the two graphs will asymptotically approach some constants \( S_{A, \text{final}} \) and \( S_{B, \text{final}} \).
          \end{answer}
    \item After the systems reach equilibrium, the multiplicity function of \( A \) is multiplied by a constant \( C \),
          \begin{equation}
          \label{eq:6}
          g_A(U_A) \rightarrow g^{\prime}_A(U_A) = C g_A(U_A).
          \end{equation}
          Write down the new entropy \(\sigma^{\prime}_A\) and show that the temperature of system \( A \) is unchanged.
          \begin{proof}
            We have that
            \begin{equation}
            \label{eq:7}
            \sigma^{\prime}_A = \log(g^{\prime}_A(U_A)) = \log C + \log g_A(U_A) = \log C + \sigma_A.
            \end{equation}

            Then, at equilibrium, the temperature satisfies
            \begin{align}
            \label{eq:8}
              \tau^{\prime} &= \left( \frac{\partial \sigma^{\prime}_A}{\partial U}  \right)^{-1} = \left( \frac{\partial }{\partial U} \log C + \frac{\partial }{\partial U} \log g_A(U_A) \right)^{-1} \\
              &= \left( \frac{\partial }{\partial U}  \log g_A(U_A) \right) = \tau.
            \end{align}
          \end{proof}
    \item What is the resulting change in the energy of system \( A \)? Of system \( B \)?
          \begin{answer}
            Neither change.
            The multiplicative factor shifts the graph of entropy against entropy up/down, but doesn't impact the slope.
          \end{answer}
  \end{enumerate}
\end{problem}

\begin{problem}{3}
  Suppose that the nuclei of atoms in a particular material each have spin \( \frac{3}{2} \).
  According to quantum mechanics, each nucleus can be in one of four states, labeled
  \begin{equation}
  \label{eq:9}
  m = -3/2, -1/2, 1/2, 3/2.
  \end{equation}
  Due to the crystal structure of the material, the states with \( m = \pm 3/2 \) have energy \( \varepsilon \), and the states with \( m = \pm 1/2 \) have energy \( 0 \).

  \begin{enumerate}[label=(\alph*)]
          \item Find the partition function \( Z \) for one atom in this system.
          \begin{answer}
            \( Z(\tau) = \sum\limits_s \exp( \frac{-\varepsilon_{s}}{\tau} ) \), so
            \begin{equation}
            \label{eq:10}
            Z(\tau) = 2 + \exp( -\frac{\varepsilon_{3/2}}{\tau} ) + \exp( \frac{-\varepsilon_{-3/2}}{\tau} ) = \boxed{2 + 2\exp( \frac{-\varepsilon}{\tau} )}.
            \end{equation}
          \end{answer}
          \item Find the average energy \( U \) of \( N \) nuclei at a temperature \(\tau\).
          \begin{answer}
            \begin{equation}
            \label{eq:11}
            \langle U \rangle = N \sum\limits_s \varepsilon_s \mathbb{P}(\varepsilon = \varepsilon_s) = N \varepsilon (\frac{\exp(\frac{-\varepsilon}{\tau})}{Z}) +N \varepsilon (\frac{\exp(\frac{-\varepsilon}{\tau})}{Z}),
            \end{equation}
            so
            \begin{equation}
            \label{eq:13}
            \langle U \rangle = N \left( \frac{2\varepsilon\exp(\frac{-\varepsilon}{\tau})}{2 + 2\exp(\frac{-\varepsilon}{\tau})} \right) = \boxed{\frac{N\varepsilon}{\exp(\frac{\varepsilon}{\tau}) + 1}}.
            \end{equation}

          \end{answer}
    \item Find the heat capacity \( C_V \) of \( N \) nuclei at a temperature \( \tau \).
          \begin{answer}
            The heat capacity is defined as \( \frac{\partial \langle U \rangle}{\partial \tau}  \).

            Hence,
            \begin{align}
            \label{eq:14}
              C_V &= \frac{\partial }{\partial \tau} \left( \frac{N\varepsilon}{\exp(\frac{\varepsilon}{\tau}) + 1} \right) = \frac{-N\varepsilon}{(\exp(\frac{\varepsilon}{\tau}) + 1)^2} \exp(\varepsilon/\tau) \left( \frac{-\varepsilon}{\tau^2} \right) \\
              &= \boxed{\frac{N \varepsilon^2}{\tau^2} \frac{\exp( \frac{\varepsilon}{\tau} )}{(\exp(\frac{\varepsilon}{\tau}) + 1)^2}}.
            \end{align}
          \end{answer}
  \end{enumerate}
\end{problem}

\begin{problem}{4}
  Consider two of the states of a system \( S \), \( A \) and \( B \), with energies \( \varepsilon_A \) and \( \varepsilon_B \), \( \varepsilon_A > \varepsilon_B \).
  System \( S \) is in thermal equilibrium with another system at temperature \(\tau\).
  Let \( P(A) \) and \( P(B) \) be the probability that \( S \) is in state \( A  \) and \( B \) at any given moment.
  \begin{enumerate}[label=(\alph*)]
    \item Find the limit of \( P(A)/P(B) \) as \( \tau \rightarrow 0 \).
          \begin{answer}
            \begin{equation}
            \label{eq:15}
            \frac{P(A)}{P(B)} = \frac{\exp(-\varepsilon_A/\tau)}{\exp(-\varepsilon_B/\tau)} = \exp(\frac{-(\varepsilon_{A}-\varepsilon_B)}{\tau}).
            \end{equation}
            Clearly, as \(\tau \rightarrow 0\), \( \boxed{P(A)/P(B) \rightarrow 0} \).
          \end{answer}
    \item As \(\tau \rightarrow \infty\).
          \begin{answer}
            Again from \cref{eq:15}, it is clear that as \(\tau \rightarrow \infty\), the ratio \( \boxed{P(A)/P(B) \rightarrow 1} \).
          \end{answer}
    \item For \(\tau \gg \varepsilon_A - \varepsilon_B\), find \( P(A)/P(B) \) to first order.
          \begin{answer}
            If \(\tau \gg \varepsilon_A - \varepsilon_B\), then \( 1 \gg \frac{\varepsilon_A - \varepsilon_B}{\tau} \).
            We can therefore expand the function \( f(x) = \exp(-x) \) about \( 0 \).

            For \( x \ll 1 \),
            \begin{equation}
            \label{eq:16}
            e^{-x} = 1 - x + \frac{x^2}{2} - \frac{x^{3}}{3!} + \dots
            \end{equation}

            Hence we have that
            \begin{equation}
            \label{eq:17}
            \frac{P(A)}{P(B)} = \boxed{1 - \frac{\varepsilon_{A} - \varepsilon_B}{\tau} + \mathcal{O}\left((\frac{\varepsilon_{A}-\varepsilon_B}{\tau})^2\right)}.
            \end{equation}
          \end{answer}
    \item At what temperature \(\tau\) is \( P(A)/P(B) = 1/2\)?
          \begin{answer}
            \begin{align}
            \label{eq:18}
              &\exp(\frac{-(\varepsilon_{A} - \varepsilon_B)}{\tau} = \frac{1}{2}) \\
              &\Rightarrow \frac{-(\varepsilon_{A} - \varepsilon_{B})}{\tau} = \log \frac{1}{2} = -\log 2 \\
              &\Rightarrow \boxed{\tau = \frac{\varepsilon_{A} - \varepsilon_B}{\log 2}}.
            \end{align}
          \end{answer}
  \end{enumerate}
\end{problem}
\end{document}
